\documentclass[preview]{standalone}

\usepackage{amsmath}
\usepackage{amssymb}
\usepackage{stellar}
\usepackage{definitions}

\begin{document}

\id{module-generators}
\genpage

\section{Generated submodules}

\begin{snippetproposition}{generated-submodule-proposition}{Submodule generated by a subset}
    Let \(A\) be a \ring, let \(M\) be a left \(A\)-\module, and let
    \(\mathcal{X}\subseteq M\). The smallest \submodule of \(M\) containing
    \(\mathcal{X}\) is
    \[
        \langle\mathcal{X}\rangle_A
        =
        \bigcap_{\substack{N\submoduleleq M\\
        \mathcal{X}\subseteq N}}N.
    \]
    If \(\mathcal{X}=\emptyset\), then
    \(\langle\mathcal{X}\rangle_A=\{0_M\}\). Otherwise,
    \[
        \langle\mathcal{X}\rangle_A
        =
        \left\{
            a_1v_1+\cdots+a_rv_r
            \suchthat
            r\geq1,\ a_i\in A,\ v_i\in\mathcal{X}
        \right\}.
    \]
\end{snippetproposition}

\begin{snippetproof}{generated-submodule-proposition-proof}{generated-submodule-proposition}{Submodule generated by a subset}
    If \(\mathcal{X}=\emptyset\), then \(\{0_M\}\) is among the
    \submodule[submodules] containing \(\mathcal{X}\), so their intersection is
    \(\{0_M\}\).

    Suppose \(\mathcal{X}\neq\emptyset\), and denote by
    \(\widetilde N\) the set of finite linear combinations displayed in
    the statement. It contains \(\mathcal{X}\), since \(v=1_Av\), and it
    is closed under subtraction and left scalar multiplication. Thus
    \(\widetilde N\submoduleleq M\).

    Every \submodule containing \(\mathcal{X}\) contains all finite linear
    combinations of its elements, and hence contains \(\widetilde N\).
    Conversely, \(\widetilde N\) is itself one of the \submodule[submodules] in the
    intersection. Therefore
    \[
        \langle\mathcal{X}\rangle_A=\widetilde N.
    \]
\end{snippetproof}

\begin{snippetdefinition}{module-generating-families-free-module-definition}{Generating families, bases, and free modules}
    Let \(A\) be a \ring and let \(M\) be a left \(A\)-\module.
    \begin{enumerate}
        \item A \set \(\mathcal{X}\subseteq M\) is a
        \emph{generating family} if
        \[
            \langle\mathcal{X}\rangle_A=M.
        \]
        The module is \emph{finitely generated} if it admits a finite
        generating family.
        \item A \set \(\mathcal{X}\subseteq M\) is \emph{free} if,
        for all distinct \(v_1,\ldots,v_r\in\mathcal{X}\) and
        \(a_1,\ldots,a_r\in A\),
        \[
            a_1v_1+\cdots+a_rv_r=0_M
            \implies
            a_1=\cdots=a_r=0_A.
        \]
        \item A free generating family is a \emph{basis}. The module
        \(M\) is \emph{free} if it has a basis.
    \end{enumerate}
\end{snippetdefinition}

\begin{snippetexample}{polynomial-ring-free-module-example}{A polynomial ring as a free module}
    Let \(A\) be a \commutativering. The \(A\)-\module \(A[x]\) is free
    with basis
    \[
        \{1,x,x^2,\ldots\}.
    \]
    Indeed, every polynomial has a unique expression as a finite
    \(A\)-linear combination of these powers.
\end{snippetexample}

\begin{snippetexample}{matrix-natural-module-not-free-example}{A module that is not free}
    Let \(K\) be a \field,
    \[
        R=\matrices_2(K),
        \qquad
        M=K^2,
    \]
    and let \(R\) act on \(M\) by matrix-vector multiplication. The
    standard vectors \(e_1,e_2\) generate \(M\), but this family is not
    free because
    \[
        \begin{pmatrix}
            0&1\\
            0&0
        \end{pmatrix}e_1=0_M
    \]
    with a nonzero coefficient.

    In fact, \(M\) has no \(R\)-basis. A nonzero \freemodule over \(R\)
    contains the free cyclic \submodule generated by any basis element,
    which is isomorphic to \(R\). This is impossible here because
    \[
        \lineardim_K R=4>2=\lineardim_K M.
    \]
\end{snippetexample}

\begin{snippetnote}{one-generated-module-need-not-be-free}{A one-generated module need not be free}
    Let \(K\) be a \field, and let
    \[
        A=K[x],
        \qquad
        M=A/(x^2-1).
    \]
    The element \(1+(x^2-1)\) generates \(M\), but \(M\) is not free:
    the nonzero scalar \(x^2-1\) annihilates every element of \(M\).
    Thus an \(A\)-\module generated by one element need not be
    isomorphic to \(A\).
\end{snippetnote}

\section{Universal property}

\begin{snippetproposition}{finite-free-module-universal-property}{Universal property of a finite free module}
    Let \(A\) be a \ring, let \(M_1,M_2\) be left
    \(A\)-\module[modules], and let
    \(M_1\) be free with finite basis
    \[
        \mathcal{X}=\{v_1,\ldots,v_r\}.
    \]
    For every \(w_1,\ldots,w_r\in M_2\), there exists a unique
    \(A\)-\modulehomomorphism
    \[
        \varphi\colon M_1\fromto M_2
    \]
    satisfying \(\varphi(v_i)=w_i\) for every \(i\).
\end{snippetproposition}

\begin{snippetproof}{finite-free-module-universal-property-proof}{finite-free-module-universal-property}{Universal property of a finite free module}
    Every \(v\in M_1\) has a unique expression
    \[
        v=a_1v_1+\cdots+a_rv_r.
    \]
    Define
    \[
        \varphi(v)=a_1w_1+\cdots+a_rw_r.
    \]
    Uniqueness of the coordinates makes this definition well-defined.
    The coordinate formulas immediately give
    \[
        \varphi(v-v')=\varphi(v)-\varphi(v'),
        \qquad
        \varphi(av)=a\varphi(v),
    \]
    so \(\varphi\) is a \modulehomomorphism, and
    \(\varphi(v_i)=w_i\).

    If \(\psi\) is another such homomorphism, then
    \[
        \psi\left(\sum_{i=1}^r a_iv_i\right)
        =
        \sum_{i=1}^r a_i\psi(v_i)
        =
        \sum_{i=1}^r a_iw_i
        =
        \varphi\left(\sum_{i=1}^r a_iv_i\right).
    \]
    Hence \(\psi=\varphi\).
\end{snippetproof}

\section{Direct sums and finite free modules}

\begin{snippetdefinition}{module-direct-sum-definition}{Direct sum of modules}
    Let \(A\) be a \ring and let \(M_1,\ldots,M_r\) be left
    \(A\)-\module[modules]. Their finite \emph{direct sum} is
    \[
        M_1\oplus\cdots\oplus M_r
        =
        \{(v_1,\ldots,v_r)\suchthat v_i\in M_i\},
    \]
    with componentwise addition and scalar multiplication:
    \[
        (v_i)_{i=1}^r+(w_i)_{i=1}^r=(v_i+w_i)_{i=1}^r,
        \qquad
        a(v_i)_{i=1}^r=(av_i)_{i=1}^r.
    \]
\end{snippetdefinition}

\begin{snippetexample}{finite-rank-standard-free-module-example}{The standard free module}
    Let \(A\) be a \ring. The \module
    \[
        A^r=A\oplus\cdots\oplus A
    \]
    is free with standard basis
    \[
        e_1=(1_A,0_A,\ldots,0_A),
        \quad\ldots,\quad
        e_r=(0_A,\ldots,0_A,1_A).
    \]
    Every element has the unique expression
    \[
        (a_1,\ldots,a_r)=a_1e_1+\cdots+a_re_r.
    \]
\end{snippetexample}

\begin{snippetcorollary}{finite-free-module-standard-form}{Standard form of a finite free module}
    Let \(A\) be a \ring. If a left \(A\)-\module \(M\) has a basis with
    \(r\) elements, then
    \[
        M\moduleisomorphic A^r.
    \]
\end{snippetcorollary}

\begin{snippetproof}{finite-free-module-standard-form-proof}{finite-free-module-standard-form}{Standard form of a finite free module}
    Let \(v_1,\ldots,v_r\) be a basis of \(M\), and let
    \(e_1,\ldots,e_r\) be the standard basis of \(A^r\). The universal
    property gives homomorphisms
    \[
        \varphi\colon M\fromto A^r,
        \qquad
        \psi\colon A^r\fromto M
    \]
    with \(\varphi(v_i)=e_i\) and \(\psi(e_i)=v_i\). Both
    \(\psi\circ\varphi\) and \(\varphi\circ\psi\) fix their respective
    bases, so uniqueness in the universal property makes them identity
    maps. Thus \(\varphi\) and \(\psi\) are inverse isomorphisms.
\end{snippetproof}

\begin{snippetcorollary}{finitely-generated-module-free-quotient}{Finitely generated modules are quotients of finite free modules}
    Let \(A\) be a \ring and let \(M\) be a finitely generated left
    \(A\)-\module. There exist \(r\geq0\) and
    \(N\submoduleleq A^r\) such that
    \[
        M\moduleisomorphic A^r/N.
    \]
\end{snippetcorollary}

\begin{snippetproof}{finitely-generated-module-free-quotient-proof}{finitely-generated-module-free-quotient}{Finitely generated modules are quotients of finite free modules}
    Let \(v_1,\ldots,v_r\) generate \(M\). The universal property gives
    a homomorphism
    \[
        \varphi\colon A^r\fromto M,
        \qquad
        \varphi(e_i)=v_i.
    \]
    It is surjective because
    \[
        a_1v_1+\cdots+a_rv_r
        =
        \varphi(a_1e_1+\cdots+a_re_r).
    \]
    The first isomorphism theorem therefore gives
    \[
        M
        \moduleisomorphic
        A^r/\moduleker\varphi.
    \]
    Take \(N=\moduleker\varphi\).
\end{snippetproof}

\begin{snippetproposition}{finite-free-module-rank-invariance}{Invariant basis number for commutative rings}
    Let \(A\) be a nonzero \commutativering and let \(M\) be a finite free
    \(A\)-\module. If \(M\) has bases with respectively \(r\) and \(s\)
    elements, then \(r=s\).
\end{snippetproposition}

\begin{snippetproof}{finite-free-module-rank-invariance-proof}{finite-free-module-rank-invariance}{Invariant basis number for commutative rings}
    Choose a \maximalideal \(\mathfrak{m}\idealin A\), and put
    \(k=A/\mathfrak{m}\). Then \(k\) is a \field and
    \(M/\mathfrak{m}M\) is naturally a \(k\)-vector space.

    The two bases give module isomorphisms
    \(M\moduleisomorphic A^r\) and
    \(M\moduleisomorphic A^s\). Passing to the quotient by
    \(\mathfrak{m}\) gives
    \[
        M/\mathfrak{m}M
        \moduleisomorphic
        (A/\mathfrak{m})^r=k^r
    \]
    and
    \[
        M/\mathfrak{m}M
        \moduleisomorphic
        (A/\mathfrak{m})^s=k^s.
    \]
    Finite-dimensional vector spaces have invariant dimension, so
    \(r=s\).
\end{snippetproof}

\begin{snippetdefinition}{finite-free-module-rank-definition}{Rank of a finite free module}
    Let \(A\) be a nonzero \commutativering and let \(M\) be a finite free
    \(A\)-\module. The common cardinality of all finite bases of \(M\) is
    the \emph{rank} of \(M\), denoted
    \[
        \freemodulerank_A M.
    \]
\end{snippetdefinition}

\section{Abelian groups as modules}

\begin{snippetexample}{abelian-groups-as-integer-modules-example}{Abelian groups as integer modules}
    Every \abeliangroup \(G\), written additively, is an
    \(\integers\)-\module under
    \[
        n\cdot g=
        \begin{cases}
            \underbrace{g+\cdots+g}_{n\text{ times}},&n>0,\\
            \underbrace{(-g)+\cdots+(-g)}_{-n\text{ times}},&n<0,\\
            0_G,&n=0.
        \end{cases}
    \]
    Conversely, the additive group of every \(\integers\)-\module is
    abelian.
\end{snippetexample}

\begin{snippetnote}{integer-module-multiplicative-notation}{Multiplicative notation}
    If the \abeliangroup \(G\) is written multiplicatively, the
    \(\integers\)-action is
    \[
        n\cdot g=
        \begin{cases}
            g^n,&n>0,\\
            (g^{-1})^{-n},&n<0,\\
            1_G,&n=0.
        \end{cases}
    \]
\end{snippetnote}

\end{document}
